\item \textbf{Problema de repaso: Salto de bungee.}
Una saltadora de bungee de $65.00\text{ kg}$ salta de un puente con una cuerda ligera amarrada a ella y al puente. La longitud no estirada de la cuerda es de $11.0\text{ m}$. La saltadora alcanza el fondo de su movimiento $36.0\text{ m}$ abajo del puente antes de rebotar de regreso. Queremos encontrar el intervalo de tiempo entre que deja el puente y llega al fondo de su movimiento. Su movimiento total se puede separar en una caída libre de $11.0\text{ m}$ y una sección de $25.0\text{ m}$ de oscilación armónica simple.
\begin{enumerate}
    \item Para la parte de caída libre, ¿cuál es el análisis de modelo adecuado para describir su movimiento?
    \item ¿Durante qué intervalo de tiempo está en caída libre?
    \item Para la parte de oscilación armónica simple de la caída, ¿es el sistema de la saltadora, la cuerda elástica y la Tierra aislado o no aislado?
    \item A partir de su respuesta en el inciso (3), encuentre la constante elástica de la cuerda.
    \item ¿Cuál es la ubicación del punto de equilibrio donde la fuerza del resorte equilibra la fuerza gravitacional ejercida sobre la saltadora?
    \item ¿Cuál es la frecuencia angular de la oscilación?
    \item ¿Qué intervalo de tiempo se requiere para que la cuerda se estire $25.0\text{ m}$?
    \item ¿Cuál es el intervalo de tiempo total para un salto de $36.0\text{ m}$?
\end{enumerate}